% ============================================================================
\section{Dynamic-Halting Protocol}
\label{app:halt-protocol}

This appendix gives the dynamic-halting protocol summarized in Section~\ref{sec:depth} and Table~\ref{tab:depth}.

\begin{table}[h]
\centering
\small
\caption{\textbf{High halting speedup can reflect $K$-invariance rather than useful adaptive computation.} Calibrated logit-margin halting on per-loop-loss models, mean over 3 seeds, timed slice harness. Thresholds are tuned on a held-out calibration slice to keep dynamic perplexity (PPL) within $1\%$ of fixed $K=4$; the table reports a separate timed slice, so small excesses over the budget reflect calibration generalization.}
\label{tab:halting-main}
\begin{tabular}{llcccc}
\toprule
\textbf{Readout} & \textbf{Model} & \textbf{Fixed $K=4$ PPL} & \textbf{Dynamic PPL} & \textbf{Avg.\ loops} & \textbf{Speedup} \\
\midrule
RMSNorm        & 44M  & 6.01 & 6.01 & \textbf{1.00} & 3.39$\times$ \\
RMSNorm        & 129M & 5.32 & 5.34 & \textbf{1.00} & 3.77$\times$ \\
Raw            & 44M  & 5.29 & 5.34 & 2.16 & 1.55$\times$ \\
Raw            & 129M & 4.86 & 4.91 & 1.76 & 2.09$\times$ \\
Final-only norm & 44M  & 5.37 & 5.41 & 1.78 & 1.90$\times$ \\
Final-only norm & 129M & 4.85 & 4.90 & 1.88 & 1.97$\times$ \\
Norm penalty   & 44M  & 5.36 & 5.41 & 2.60 & 1.29$\times$ \\
Norm penalty   & 129M & 4.85 & 4.90 & 2.56 & 1.43$\times$ \\
\bottomrule
\end{tabular}
\end{table}

\paragraph{Halting decision: per-sequence, not per-token.}
At inference time, the model runs the shared block $K_{\max}$ times.
At each loop $k$, for a sequence of length $L$, we compute a sequence-level confidence score $c_k$ as the mean over tokens of the top-1-vs-top-2 logit margin at that loop:
\begin{equation}
    c_k \;=\; \frac{1}{L}\sum_{t=1}^{L} \bigl(\mathrm{logits}^{(1)}_{k,t} - \mathrm{logits}^{(2)}_{k,t}\bigr),
\end{equation}
where the two logits are the largest and second-largest at position $t$.
A sequence halts at the smallest $k$ for which $c_k$ exceeds a fixed threshold $\tau$.
A sequence that never crosses $\tau$ runs to $K_{\max}$.
We use the mean over tokens rather than the minimum because we found the minimum to be brittle to single-token uncertainty in long sequences---a single hard token at the end forces the whole sequence to spend $K_{\max}$ loops.
Per-token early exit is a different problem (it interacts non-trivially with autoregressive decoding and key-value (KV) caching) and is outside the scope of this paper.

\paragraph{Threshold calibration.}
Given a fixed $K_{\max}{=}4$ and a target relative-PPL budget $\rho{=}1.01$ (within $1\%$ of the fixed-$K_{\max}$ PPL), we sweep $\tau$ on a held-out calibration set of $160$ sequences and pick the smallest $\tau$ such that the resulting dynamic PPL is within $\rho$ of the fixed-$K_{\max}$ PPL on the same sequences.
The calibration is done per checkpoint, so $\tau$ varies across conditions and seeds.
Sweeping $\tau$ uses cached per-loop logits from a single forward pass at $K_{\max}$, so calibration is cheap.

\paragraph{Throughput benchmark.}
After calibration we run a separate timed benchmark on a different 160-sequence slice with the same protocol.
The forward pass uses dynamic batching: at each loop $k$, only sequences that have not yet halted continue to the next loop.
Sequences that halt at loop $k$ leave the batch.
This gives the wall-clock benefit of stopping early.
The reported throughput in the dynamic-halting benchmark is total tokens processed (across all sequences and all positions) divided by total benchmark wall-clock time, including the cost of removing halted sequences from the batch.
Timing setup: single GPU, batch size 8, sequence length 256, bfloat16 (bf16) inference, no KV caching (we evaluate teacher-forced language-model (LM) scoring rather than autoregressive generation).
The fixed-$K_{\max}$ comparison uses the same batch geometry on the same slice.
The reported $\mathrm{Speedup}=\mathrm{tokens/s}_{\mathrm{dynamic}}/\mathrm{tokens/s}_{\mathrm{fixed}\,K_{\max}}$ is therefore mostly attributable to the difference in average loop count, not to batching artifacts.

\paragraph{Caveats.}
This protocol measures speedup in teacher-forced LM scoring, which is the standard variable-depth evaluation in prior work but is not the same as autoregressive generation throughput.
A real-time generation pipeline would interact with KV caching, decoding sampling, and per-step batch-shape changes in ways that we do not measure here.
The protocol is also calibrated on a relatively short held-out slice; with longer sequences or distribution shift, $\tau$ would need recalibration.
We use the same protocol across all conditions, so within-row comparisons are clean even if absolute speedups would differ in a production setting.
